\documentclass[11pt]{article}
\usepackage[utf8]{inputenc}
\usepackage[T1]{fontenc}
\usepackage{amsmath}
\usepackage{geometry}
\geometry{margin=2.5cm}

\title{SCCI-2026: A Local Geometric Diagnostic for Bridge Detection}
\author{EHK-SCCI / CATALYST Working Document}
\date{24 November 2026}

\begin{document}

\maketitle

\section{Scope}

This working document describes a local reproducibility and feasibility
artefact for the SCCI-2026 Layer 1 geometric diagnostic workflow.

The package demonstrates exact Wasserstein-1 transport, radius-dependent
edge inclusion, and Ollivier-Ricci curvature computation on a canonical
two-triangle graph connected by a bridge edge.

This document does not validate the complete EHK-SCCI/CATALYST system,
Implication (A), Implication (B), the target equivalence, spectral dynamics,
the De Sitter attractor, or any physical, economic, or public-policy claim.

\section{Canonical graph}

The canonical graph consists of two triangles connected by the bridge edge
$(2,3)$:

\[
1-2-4, \qquad 3-5-6, \qquad 2-3.
\]

The bridge edge connects the two dense local communities. The transport pair
$(1,5)$ has geodesic distance $3$.

\section{Radius policy}

With $\alpha=0.5$, the validator applies two radius policies:

\begin{itemize}
  \item At $r_{\max}=2$, the transport pair distance $3$ exceeds the radius.
  Therefore, the bridge edge $(2,3)$ is marked \texttt{EXCLUDED}.
  \item At $r_{\max}=3$, the bridge edge is included and the exact
  Wasserstein-1 transport is resolved.
\end{itemize}

\section{Curvature result}

The exact Wasserstein-1 transport between the neighbourhood measures of the
bridge endpoints gives:

\[
W_1(\mu_2,\mu_3)=\frac{4}{3}.
\]

Since the bridge distance is $1$, the Ollivier-Ricci curvature is:

\[
\kappa_{0.5}(2,3)
=
1-\frac{W_1(\mu_2,\mu_3)}{d(2,3)}
=
1-\frac{4}{3}
=
-\frac{1}{3}.
\]

Numerically, the validator reports:

\[
\kappa_{0.5}(2,3)=-0.3333333333333337.
\]

This value is consistent with $-1/3$ within the specified tolerance of
$10^{-9}$.

\section{Interpretation}

Dense local communities tend to exhibit non-negative curvature, whereas the
bridge edge acts as a structural bottleneck and receives negative curvature.
This behaviour supports the use of SCCI-style geometric diagnostics as an
interpretable signal for structural analysis.

The result is a local numerical demonstration. It is not a proof of a
universal anomaly-detection property, nor a validation of the complete
EHK-SCCI/CATALYST framework.

\section{Reproducibility}

Run the following commands from the repository root:

\begin{verbatim}
python3 -m unittest discover -s tests -v
python3 scripts/validate_scci.py
sha256sum --check --strict SHA256SUMS
\end{verbatim}

The validator writes:

\begin{verbatim}
results/canonical_bridge_audit_r3.json
results/canonical_bridge_audit_r2_excluded.json
\end{verbatim}

\section{Governance}

The workflow follows a fail-closed policy: missing inputs, inconsistent
parameters, or failed checks produce a failure state rather than a silent
pass.

This approach is aligned methodologically with the Govern, Map, Measure and
Manage functions of the NIST AI Risk Management Framework 1.0.

\end{document}
